\section{Additional remarks}\label{remarks-sec}

The proof of Theorem \ref{main-dhl}(xii) may be modified to establish the following variant:

\begin{proposition}  Assume the generalized Elliott-Halberstam conjecture $\GEH[\theta]$ for all $0 < \theta < 1$.  Let $0 < \eps < 1/2$ be fixed.  Then if $x$ is a sufficiently large multiple of $6$, there exists a natural number $n$ with $\eps x \leq n \leq (1-\eps) x$ such that at least two of $n, n-2, x-n$ are prime.  Similarly if $n-2$ is replaced by $n+2$.
\end{proposition}

Note that if at least two of $n,n-2,x-n$ are prime, then either $n,n+2$ are twin primes, or else at least one of $x,x-2$ is expressible as the sum of two primes, and Theorem \ref{disj} easily follows.

\begin{proof} (Sketch) We just discuss the case of $n-2$, as the $n+2$ case is similar. Observe from the Chinese remainder theorem (and the hypothesis that $x$ is divisible by $6$) that one can find a residue class $b\ (W)$ such that $b, b-2, x-b$ are all coprime to $W$ (in particular, one has $b=1\ (6)$).   By a routine modification of the proof of Lemma \ref{crit}, it suffices to  find a non-negative weight
  function $\nu \colon \N \to \R^+$ and fixed quantities $\alpha > 0$ and $\beta_1,\beta_2,\beta_3 \geq
  0$, such that one has the
  asymptotic upper bound
$$
 \sum_{\substack{\eps x \leq n \leq (1-\eps) x\\ n = b\ (W)}} \nu(n) \leq {\mathfrak S} (\alpha+o(1)) B^{-k} \frac{(1-2\eps) x}{W},
$$
the asymptotic lower bounds
\begin{align*}
  \sum_{\substack{\eps x \leq n \leq (1-\eps) x\\ n = b\ (W)}} \nu(n) \theta(n) &\geq {\mathfrak S} (\beta_1-o(1)) B^{1-k} \frac{(1-2\eps) x}{\phi(W)} \\
  \sum_{\substack{\eps x \leq n \leq (1-\eps) x\\ n = b\ (W)}} \nu(n) \theta(n+2) &\geq {\mathfrak S} (\beta_2-o(1)) B^{1-k} \frac{(1-2\eps) x}{\phi(W)} \\
  \sum_{\substack{\eps x \leq n \leq (1-\eps) x\\ n = b\ (W)}} \nu(n) \theta(x-n) &\geq {\mathfrak S} (\beta_3-o(1)) B^{1-k} \frac{(1-2\eps) x}{\phi(W)} 
\end{align*}
and the inequality
$$ \beta_1+\beta_2+\beta_3 > 2 \alpha,$$
where ${\mathfrak S}$ is the singular series
$$ {\mathfrak S} := \prod_{p|x(x-2); p > w} \frac{p}{p-1}.$$
We select $\nu$ to be of the form
$$ \nu(n) = \left( \sum_{j=1}^J c_j \lambda_{F_{j,1}}(n) \lambda_{F_{j,2}}(n+2) \lambda_{F_{j,3}}(x-n) \right)^2 $$
for various fixed coefficients $c_1,\dots,c_J \in \R$ and fixed smooth compactly supported functions $F_{j,i}: [0,+\infty) \to \R$ with $j=1,\dots,J$ and $i=1,\dots,3$.  It is then routine\footnote{One new technical difficulty here is that some of the various moduli $[d_j,d'_j]$ arising in these arguments are not required to be coprime at primes $p > w$ dividing $x$ or $x-2$; this requires some modification to Lemma \ref{mul-asym} that ultimately leads to the appearance of the singular series ${\mathfrak S}$.  However, these modifications are quite standard, and we do not give the details here.} to verify that analogues of Theorem \ref{prime-asym} and Theorem \ref{nonprime-asym} hold for the various components of $\nu$, with the role of $x$ in the right-hand side replaced by $(1-2\eps) x$, and the claim then follows by a suitable modification of Theorem \ref{epsilon-beyond}, taking advantage of the function $F$ constructed in Theorem \ref{piece}.
\end{proof}

It is likely that the bounds in Theorem \ref{main} can be improved further by refining the sieve-theoretic methods employed in this paper, with the exception of part (xii) for which the parity problem prevents further improvement, as discussed in Section \ref{parity-sec}.  We list some possible avenues to such improvements as follows:

\begin{enumerate}
\item In Theorem \ref{mke-lower}, the bound $M_{k,\eps} > 4$ was obtained for some $\eps>0$ and $k=50$.  It is possible that $k$ could be lowered slightly, for instance to $k=49$, by further numerical computations, but we were only barely able to establish the $k=50$ bound after two weeks of computation.  However, there may be a more efficient way to solve the required variational problem (e.g. by selecting a more efficient basis than the symmetric monomial basis) that would allow one to advance in this direction; this would improve the bound $H_1 \leq 246$ slightly.
\item To reduce $k$ (and thus $H_1$) further, one could try to solve another variational problem, such as the one arising in Theorem \ref{maynard-trunc} or in Theorem \ref{epsilon-beyond}, rather than trying to lower bound $M_k$ or $M_{k,\eps}$.  It is also possible to use the more complicated versions of $\MPZ[\varpi,\delta]$ established in \cite{polymath8a} (in which the modulus $q$ is assumed to be densely divisible rather than smooth) to replace the truncated simplex appearing in Theorem \ref{maynard-trunc} with a more complicated region (such regions also appear implicitly in \cite[\S 4.5]{polymath8a}).  However, in the medium-dimensional setting $k \approx 50$, we were not able to accurately and rapidly evaluate the various integrals associated to these variational problems when applied to a suitable basis of functions.  One key difficulty here is that whereas polynomials appear to be an adequate choice of basis for the $M_k$, an analysis of the Euler-Lagrange equation reveals that one should use piecewise polynomial basis functions instead for more complicated variational problems such as the $M_{k,\eps}$ problem (as was done in the three-dimensional case in Section \ref{3d}), and these are difficult to work with in medium dimensions.  From our experience with the low $k$ problems, it looks like one should allow these piecewise polynomials to have relatively high degree on some polytopes, low degree on other polytopes, and vanish completely on yet further polytopes\footnote{In particular, the optimal choice $F$ for $M_{k,\eps}$ should vanish on the polytope $\{ (t_1,\dots,t_k) \in (1+\eps) \cdot {\mathcal R}_k: \sum_{i \neq i_0} t_i \ge 1-\eps \hbox{ for all } i_0=1,\dots,k\}$.}, but we do not have a systematic understanding of what the optimal placement of degrees should be.  
\item In Theorem \ref{epsilon-beyond}, the function $F$ was required to be supported in the simplex $\frac{k}{k-1} \cdot {\mathcal R}_k$.  However, one can consider functions $F$ supported in other regions $R$, subject to the constraint that all elements of the sumset $R+R$ lie in a region treatable by one of the cases of Theorem \ref{nonprime-asym}.  This could potentially lead to other optimization problems that lead to superior numerology, although again it appears difficult to perform efficient numerics for such problems in the medium $k$ regime $k \approx 50$.  One possibility would be to adopt a ``free boundary'' perspective, in which the support of $F$ is not fixed in advance, but is allowed to evolve by some iterative numerical scheme.
\item To improve the bounds on $H_m$ for $m=2,3,4,5$, one could seek a better lower bound on $M_k$ than the one provided by Theorem \ref{explicit}; one could also try to lower bound more complicated quantities such as $M_{k,\eps}$.
\item One could attempt to improve the range of $\varpi,\delta$ for which estimates of the form $\MPZ[\varpi,\delta]$ are known to hold, which would improve the results of Theorem \ref{main}(ii)-(vi).  For instance, we believe that the condition $600 \varpi + 180\delta < 7$ in Theorem \ref{mpz-poly} could be improved slightly to $1080 \varpi + 330 \delta < 13$ by refining the arguments in \cite{polymath8a}, but this requires a hypothesis of square root cancellation in a certain four-dimensional exponential sum over finite fields, which we have thus far been unable to establish rigorously.  Another direction to pursue would be to improve the $\delta$ parameter, or to otherwise relax the requirement of smoothness in the moduli, in order to reduce the need to pass to a truncation of the simplex ${\mathcal R}_k$, which is the primary reason why the $m=1$ results are currently unable to use the existing estimates of the form $\MPZ[\varpi,\delta]$.  Another speculative possibility is to seek $\MPZ[\varpi,\delta]$ type estimates which only control distribution for a positive proportion of smooth moduli, rather than for all moduli, and then to design a sieve $\nu$ adapted to just that proportion of moduli (cf. \cite{fouvry-invent}).  Finally, there may be a way to combine the arguments currently used to prove $\MPZ[\varpi,\delta]$ with the automorphic forms (or ``Kloostermania'') methods used to prove nontrivial equidistribution results with respect to a fixed modulus, although we do not have any ideas on how to actually achieve such a combination.
\item It is also possible that one could tighten the argument in Lemma \ref{crit}, for instance by establishing a non-trivial lower bound on the portion of the sum $\sum_n \nu(n)$ when $n+h_1,\dots,n+h_k$ are all composite, or a sufficiently strong upper bound on the pair correlations $\sum_n \theta(n+h_i) \theta(n+h_j)$ (see \cite[\S 6]{banks} for a recent implementation of this latter idea).  However, our preliminary attempts to exploit these adjustements suggested that the gain from the former idea would be exponentially small in $k$, whereas the gain from the latter would also be very slight (perhaps reducing $k$ by $O(1)$ in large $k$ regimes, e.g. $k \geq 5000$).
\item All of our sieves used are essentially of Selberg type, being the square of a divisor sum.  We have experimented with a number of non-Selberg type sieves (for instance trying to exploit the obvious positivity of $1 - \sum_{p \leq x: p|n} \frac{\log p}{\log x}$ when $n \leq x$), however none of these variants offered a numerical improvement over the Selberg sieve.  Indeed it appears that after optimizing the cutoff function $F$, the Selberg sieve is in some sense a ``local maximum'' in the space of non-negative sieve functions, and one would need a radically different sieve to obtain numerically superior results.
\item Our numerical bounds for the diameter $H(k)$ of the narrowest admissible $k$-tuple are known to be exact for $k \leq 342$, but there is scope for some slight improvement for larger values of $k$, which would lead to some improvements in the bounds on $H_m$ for $m=2,3,4,5$.  However, we believe that our bounds on $H_m$ are already fairly close (e.g. within $10\%$) of optimal, so there is only a limited amount of gain to be obtained solely from this component of the argument.
\end{enumerate}
